\chapter{Research Methodology}

This chapter presents the complete research methodology of FEDMED-X. The
methodology is designed as a systems pipeline rather than a single-model
experiment. It combines data partitioning, lightweight neural architecture
design, federated optimization, formal privacy accounting, communication
compression, and evaluation protocol design under one reproducible
framework. The guiding principle is that medical federated learning must be
assessed under joint constraints of non-IID heterogeneity, privacy, bandwidth
limitation, and edge-oriented deployability~\cite{guan2024federated,
ghosh2026advancements,zhou2025federated}. Figure~\ref{fig:framework}
illustrates the end-to-end methodological flow used in this project.

% =====================================================
% PROFESSIONAL WORKFLOW DIAGRAM
% =====================================================
\begin{figure}[!htbp]
\centering
\begin{tikzpicture}[
  node distance = 7mm and 10mm,
  every node/.style = {font=\footnotesize},
  stage/.style = {rectangle, rounded corners=3pt, draw=black!70,
                  minimum width=5.6cm, minimum height=8mm,
                  align=center, thick, drop shadow},
  data/.style   = {stage, fill=green!15},
  train/.style  = {stage, fill=orange!20},
  priv/.style   = {stage, fill=red!15},
  comm/.style   = {stage, fill=blue!12},
  serv/.style   = {stage, fill=violet!15},
  eval/.style   = {stage, fill=teal!15},
  arrow/.style  = {-{Latex[length=2mm]}, thick, black!70}
]
\node[data]  (d1) {MedMNIST PneumoniaMNIST (5{,}856 images)};
\node[data,  below=of d1] (d2) {Dirichlet Partition ($\alpha=0.5$)};
\node[data,  below=of d2] (d3) {5 Simulated Hospital Clients (Non-IID)};
\node[train, below=of d3] (t1) {MedNet-CBAM Local Training + Focal Loss ($\gamma=2.0$)};
\node[priv,  below=of t1] (p1) {DP-SGD Privacy Engine ($\epsilon=1.0$, $\delta=10^{-5}$)};
\node[comm,  below=of p1] (c1) {TopK 50\% Sparsification + INT8 Quantization ($\sim$90\% BW$\downarrow$)};
\node[serv,  below=of c1] (s1) {Server: Weighted FedProx Aggregation ($\mu=0.01$)};
\node[eval,  below=of s1] (e1) {Evaluation: Accuracy, Macro-F1, Precision, Recall};
\node[eval,  below=of e1] (e2) {Edge Benchmark: NVIDIA Jetson Nano (target 0.20\,ms/img)};

\draw[arrow] (d1) -- (d2);
\draw[arrow] (d2) -- (d3);
\draw[arrow] (d3) -- (t1);
\draw[arrow] (t1) -- (p1);
\draw[arrow] (p1) -- (c1);
\draw[arrow] (c1) -- (s1);
\draw[arrow] (s1) -- (e1);
\draw[arrow] (e1) -- (e2);

% Feedback loop
\draw[arrow, dashed, black!55]
  (s1.east) .. controls +(2.5,0) and +(2.5,0) .. (t1.east)
  node[midway, right=2pt, font=\scriptsize]{Global Model Broadcast};
\end{tikzpicture}
\caption{Proposed FEDMED-X methodological workflow showing data
partitioning, local training, privacy accounting, compression, server
aggregation, evaluation, and edge-oriented benchmarking. The dashed arrow
represents the recurring global-model broadcast at every federated round.}
\label{fig:framework}
\end{figure}

\noindent\emph{Figure Explanation:} Figure~\ref{fig:framework} depicts the
six functional layers of FEDMED-X: (i)~data layer, (ii)~local training layer,
(iii)~privacy layer, (iv)~compression layer, (v)~aggregation layer, and
(vi)~evaluation layer. Each layer corresponds to a toggleable module in the
implementation, allowing the ablation configurations A1--A5 described in
Section~3.3.7 to reuse the same infrastructure without modification.
\FloatBarrier

% -------------------------------------------------
\section{Proposed Framework}
FEDMED-X is formulated as a multi-hospital federated learning framework with
five simulated clients. In every communication round, each client receives
the current global model, performs local training on its private partition,
optionally applies differentially private optimization, compresses the
outgoing update, and sends the update to the server. The server aggregates
client updates by sample-size-weighted averaging and broadcasts the new
global model. The framework is deliberately modular so that each design
choice can be isolated in ablation experiments:
\begin{enumerate}[leftmargin=*]
  \item \textbf{Baseline aggregation path:} FedAvg without privacy or
        compression;
  \item \textbf{Non-IID robust path:} FedProx with proximal regularization;
  \item \textbf{Privacy path:} FedProx + DP-SGD;
  \item \textbf{Compression path:} FedProx + TopK/INT8;
  \item \textbf{Full systems path:} FedProx + DP-SGD + TopK/INT8 (FEDMED-X).
\end{enumerate}
This factorial structure is essential for scientific interpretability. It
allows the project to measure not only final accuracy, but also the marginal
effect of privacy and compression under the same data split, model family,
and training protocol~\cite{zhou2025federated}.

\subsection{Design Principles}
The proposed framework follows four engineering principles:
\begin{itemize}[leftmargin=*]
  \item \textbf{Privacy by design:} raw images never leave the client
        boundary;
  \item \textbf{Robustness to heterogeneity:} local optimization is
        constrained to reduce harmful drift;
  \item \textbf{Communication realism:} update payloads are compressed
        before server transmission;
  \item \textbf{Edge eligibility:} the backbone remains lightweight enough
        for resource-constrained deployment targets such as NVIDIA Jetson
        Nano class devices~\cite{shabarinath2025efficacy}.
\end{itemize}

% -------------------------------------------------
\section{Data Set Analysis / Software Project Requirement}

\subsection{Dataset Selection}
The primary dataset used in this study is PneumoniaMNIST from MedMNIST
v2~\cite{yang2023medmnist}. MedMNIST is selected because it is a
standardized, reproducible biomedical benchmark that supports controlled
methodological comparison. PneumoniaMNIST is a binary classification task
with labels:
\begin{itemize}[leftmargin=*]
  \item Class 0: Normal
  \item Class 1: Pneumonia
\end{itemize}
All images are grayscale with resolution $28 \times 28$.

\begin{table}[!htbp]
\centering
\caption{Dataset summary for PneumoniaMNIST.}
\label{tab:dataset}
\small
\begin{tabular}{@{}lccc@{}}
\toprule
\textbf{Split} & \textbf{Description} & \textbf{Samples} & \textbf{Classes} \\
\midrule
Train      & Client partitioning source  & 4{,}708 & 2 \\
Validation & Model monitoring split      & 524     & 2 \\
Test       & Final held-out evaluation   & 624     & 2 \\
\midrule
Total      & Full PneumoniaMNIST task    & 5{,}856 & 2 \\
\bottomrule
\end{tabular}
\end{table}

\noindent\emph{Table Explanation:} Table~\ref{tab:dataset} shows the
official PneumoniaMNIST split used throughout FEDMED-X. The training split
is further redistributed across five clients using Dirichlet sampling.

\begin{table}[!htbp]
\centering
\caption{Class definition used in FEDMED-X.}
\label{tab:classes}
\small
\begin{tabular}{@{}cl@{}}
\toprule
\textbf{Label} & \textbf{Clinical Meaning} \\
\midrule
0 & Normal chest X-ray \\
1 & Pneumonia-positive chest X-ray \\
\bottomrule
\end{tabular}
\end{table}

\noindent\emph{Table Explanation:} Table~\ref{tab:classes} defines the two
diagnostic labels used for both training and evaluation.
\FloatBarrier

\subsection{Non-IID Client Partitioning}
To simulate multi-hospital heterogeneity, the training set is partitioned
across five clients using a Dirichlet distribution with concentration
parameter $\alpha = 0.5$. Lower $\alpha$ produces stronger label skew, which
better reflects realistic institutional
imbalance~\cite{riyadi2026comparative,alekseenko2024distance}. Each client
$k$ receives a private index subset $\mathcal{D}_k$ such that:
\begin{equation}
\bigcup_{k=1}^{5} \mathcal{D}_k = \mathcal{D}_{\text{train}}, \qquad
\mathcal{D}_i \cap \mathcal{D}_j = \emptyset \quad (i \neq j).
\end{equation}
This creates uneven local class distributions and unequal local dataset
sizes, which stress-test federated optimization.

\subsection{Software and System Requirements}
Table~\ref{tab:software} summarizes the implementation stack.

\begin{table}[!htbp]
\centering
\caption{Software and environment requirements.}
\label{tab:software}
\small
\begin{tabular}{@{}ll@{}}
\toprule
\textbf{Component} & \textbf{Specification} \\
\midrule
Language              & Python 3.x \\
Deep learning         & PyTorch \\
Privacy engine        & Opacus (DP-SGD, R\'enyi accounting) \\
Dataset API           & MedMNIST v2 \\
Training hardware     & Google Colab NVIDIA T4 GPU \\
Target edge device    & NVIDIA Jetson Nano 4\,GB class \\
Experiment tracking   & CSV logs + saved model checkpoints \\
Report toolchain      & VS Code + LaTeX Workshop \\
\bottomrule
\end{tabular}
\end{table}

\noindent\emph{Table Explanation:} Table~\ref{tab:software} lists the tools
and hardware used across training, privacy accounting, and reproducibility.
\FloatBarrier

\subsection{Functional Requirements}
From a software engineering perspective, the system must satisfy:
\begin{enumerate}[leftmargin=*]
  \item load and preprocess PneumoniaMNIST images;
  \item partition training data into five non-IID clients;
  \item support FedAvg and FedProx training modes;
  \item optionally enable DP-SGD with configurable $\epsilon$;
  \item optionally compress updates using TopK and INT8 simulation;
  \item aggregate client updates on the server;
  \item evaluate Accuracy, Precision, Recall, and Macro-F1;
  \item export ablation results for report-level analysis.
\end{enumerate}

\subsection{Non-Functional Requirements}
\begin{itemize}[leftmargin=*]
  \item \textbf{Reproducibility:} fixed seeds and identical protocol across
        ablation runs;
  \item \textbf{Modularity:} privacy and compression can be toggled
        independently;
  \item \textbf{Efficiency:} lightweight backbone for constrained
        deployment;
  \item \textbf{Auditability:} quantitative logs for every experimental
        configuration.
\end{itemize}

% -------------------------------------------------
\section{Algorithm / Model Analysis}

\subsection{MedNet-CBAM Backbone}
The classification backbone is a custom lightweight CNN named MedNet,
augmented with CBAM attention~\cite{woo2018cbam}. The network is
intentionally compact to remain eligible for edge deployment. Important
implementation choices include:
\begin{itemize}[leftmargin=*]
  \item GroupNorm instead of BatchNorm for DP compatibility;
  \item non-inplace ReLU activations for stable autograd under Opacus;
  \item CBAM blocks after early convolutional stages to emphasize
        informative regions.
\end{itemize}
A simplified forward path is:
$x \rightarrow \text{ConvBlock1} + \text{CBAM} \rightarrow
\text{ConvBlock2} + \text{CBAM} \rightarrow \text{ConvBlock3}
\rightarrow \text{Classifier}$.

\subsection{Local Objective and Focal Loss}
Because medical class distributions can be imbalanced, local training uses
Focal Loss~\cite{lin2017focal}:
\begin{equation}
\mathcal{L}_{\text{focal}} = -\alpha_t (1 - p_t)^{\gamma} \log(p_t),
\end{equation}
with $\gamma = 2.0$. This down-weights easy examples and focuses learning on
harder samples.

\subsection{FedAvg Baseline}
In FedAvg, client $k$ minimizes its local empirical loss and returns updated
parameters $w_k$. The server aggregates~\cite{mcmahan2017communication}:
\begin{equation}
w^{t+1} = \sum_{k=1}^{K} \frac{n_k}{n} \, w_k^{t+1},
\end{equation}
where $n_k = |\mathcal{D}_k|$ and $n = \sum_k n_k$. FedAvg is effective
under mild heterogeneity but can degrade under strong non-IID skew.

\subsection{FedProx Optimization}
FedProx extends local training with a proximal
penalty~\cite{li2020fedprox,yuan2022convergence}:
\begin{equation}
\mathcal{L}_{\text{FedProx}}(w; w^t) = \mathcal{L}_{\text{local}}(w)
+ \frac{\mu}{2} \| w - w^t \|_2^2,
\end{equation}
where $\mu = 0.01$ in this project. The proximal term limits client drift
and stabilizes aggregation under heterogeneous partitions. This choice is
supported by comparative medical FL evidence on MedMNIST-style
tasks~\cite{riyadi2026comparative}.

\subsection{DP-SGD Privacy Mechanism}
When privacy mode is enabled, local optimization uses
DP-SGD~\cite{abadi2016deep}. Per-sample gradients are clipped to a maximum
norm $C$ and Gaussian noise is added before parameter update. The
implemented target budget is
\begin{equation}
\epsilon = 1.0, \quad \delta = 1 \times 10^{-5}.
\end{equation}
Privacy accounting is handled through Opacus with R\'enyi Differential
Privacy~\cite{mironov2017renyi}. This provides a formal interpretation of
privacy loss rather than an informal claim of
confidentiality~\cite{ziller2021differentially,hou2023improving}.

\subsection{TopK + INT8 Compression}
To reduce communication cost, outgoing updates are compressed in two stages:
\begin{enumerate}[leftmargin=*]
  \item \textbf{TopK sparsification (50\%):} only the highest-magnitude half
        of coefficients are retained~\cite{habib2025cflsparsemed};
  \item \textbf{INT8 quantization simulation:} retained values are quantized
        to 8-bit scale and dequantized for aggregation
        compatibility~\cite{asad2024fesec,sun2025medfd}.
\end{enumerate}
The combined scheme is designed to approximate substantial payload
reduction. In this project, the full compressed setting is associated with a
90\% bandwidth-saving target relative to dense full-precision update
exchange.

\subsection{Server Aggregation and Evaluation Protocol}
After local training and optional compression, the server performs weighted
aggregation and evaluates the global model on the held-out test set. The
main metrics are Accuracy, Macro-F1, Precision, and Recall. The
methodological evaluation is organized as a five-run factorial ablation:

\begin{table}[!htbp]
\centering
\caption{Factorial ablation configurations used in FEDMED-X.}
\label{tab:ablation}
\small
\begin{tabular}{@{}clccc@{}}
\toprule
\textbf{Run} & \textbf{Configuration} & \textbf{FedProx} & \textbf{DP-SGD} & \textbf{TopK/INT8} \\
\midrule
A1 & FedAvg baseline    & No  & No  & No  \\
A2 & FedProx baseline   & Yes & No  & No  \\
A3 & Privacy path       & Yes & Yes & No  \\
A4 & Compression path   & Yes & No  & Yes \\
A5 & Full FEDMED-X      & Yes & Yes & Yes \\
\bottomrule
\end{tabular}
\end{table}

\noindent\emph{Table Explanation:} Table~\ref{tab:ablation} defines the five
controlled ablation runs. This protocol is the methodological core of the
project because it enables causal interpretation of component effects under
one shared experimental family.
\FloatBarrier

% -------------------------------------------------
\section{Explanation of Different Modules}
FEDMED-X is implemented as a set of interdependent modules. Each module
encapsulates one responsibility in the pipeline.

\subsection{Data Module}
The data module performs dataset download, tensor transformation,
normalization, and Dirichlet partitioning. Normalization uses mean and
standard deviation of 0.5 for grayscale inputs. Outputs include: client-wise
training dataloaders, a global validation loader, and a global test loader.

\subsection{Model Module}
The model module defines MedNet-CBAM and exposes a common interface for
training and inference. GroupNorm and non-inplace activations are mandatory
design decisions for DP compatibility and stable backward computation under
privacy engines.

\subsection{Federated Training Module}
This module coordinates communication rounds. For each round it:
(i)~broadcasts global weights; (ii)~trains selected clients locally;
(iii)~collects updates; (iv)~aggregates weights; (v)~evaluates global
performance.

\subsection{Privacy Module}
The privacy module wraps local training with Opacus DP-SGD when enabled. It
controls clipping, noise assignment, and target $\epsilon$. In full-stack
mode, privacy is applied \emph{before} update compression to preserve the
declared privacy budget.

\subsection{Compression Module}
The compression module applies TopK masking and INT8 quantization simulation
on model updates. It can be disabled for baseline runs to measure
compression-only effects.

\subsection{Evaluation and Logging Module}
The evaluation module computes metrics and stores experiment-level
summaries, including accuracy, F1-score, precision, recall, privacy mode,
compression mode, and runtime. These logs are used in later chapters for
result analysis and discussion.

\subsection{Edge-Oriented Deployment Module}
Although training is performed in a GPU cloud environment, the methodology
includes an edge-oriented deployment perspective. The model footprint and
inference-oriented design choices are aligned with NVIDIA Jetson Nano class
constraints~\cite{song2025edgecloud,shabarinath2025efficacy,memon2024edge}.
The intended inference target characteristics used in this project are
ultra-low per-image latency and high throughput on constrained hardware.

\begin{figure}[!htbp]
\centering
\begin{tikzpicture}[
  node distance = 8mm and 12mm,
  every node/.style = {font=\footnotesize},
  mod/.style = {rectangle, rounded corners=3pt, draw=black!70,
                minimum width=3.2cm, minimum height=8mm,
                align=center, thick, fill=blue!10, drop shadow},
  arrow/.style = {-{Latex[length=2mm]}, thick, black!70}
]
\node[mod] (dm) {Data Module};
\node[mod, right=of dm] (mm) {Model Module};
\node[mod, right=of mm] (ft) {FL Trainer};

\node[mod, below=of dm, fill=red!12]     (pm) {Privacy Module};
\node[mod, below=of mm, fill=blue!12]    (cm) {Compression Module};
\node[mod, below=of ft, fill=violet!12]  (ag) {Aggregation};

\node[mod, below=of cm, fill=teal!12]    (ev) {Evaluation \& Logging};
\node[mod, right=of ev, fill=orange!15]  (edge) {Edge Target Profile};

\draw[arrow] (dm) -- (mm);
\draw[arrow] (mm) -- (ft);
\draw[arrow] (ft) -- (pm);
\draw[arrow] (ft) -- (cm);
\draw[arrow] (ft) -- (ag);
\draw[arrow] (pm) -- (cm);
\draw[arrow] (cm) -- (ag);
\draw[arrow] (ag) -- (ev);
\draw[arrow] (ev) -- (edge);
\end{tikzpicture}
\caption{Module-level architecture of the FEDMED-X software pipeline.}
\label{fig:modules}
\end{figure}

\noindent\emph{Figure Explanation:} Figure~\ref{fig:modules} shows the
module-level dependency structure of the implementation. Each block maps
directly to a source file in the codebase and can be toggled independently
during ablation.
\FloatBarrier

\nocite{*}

\chaptersummary{This chapter presented the methodology of FEDMED-X as a
complete systems pipeline. PneumoniaMNIST from MedMNIST v2 is partitioned
into five Dirichlet non-IID hospital clients. A lightweight MedNet-CBAM
model is trained under FedAvg and FedProx modes, with optional DP-SGD
privacy and TopK/INT8 compression. The methodology emphasizes modular design
and factorial ablation so that non-IID robustness, privacy cost, and
communication reduction can be measured under one controlled protocol. The
next chapter reports implementation details and experimental results
obtained from this methodology.}